\GalSourcePageBoundary{065}{L0064}{36}{36}

la permutation d'où l'on sera parti. Donc, si dans un pareil
groupe on a les substitutions $S$ et $T$, on est sûr d'avoir la substitution $ST$.

Telles sont les définitions que nous avons cru devoir rappeler.

\medskip
\noindent\textsc{Lemme I.} --- \emph{Une équation irréductible ne peut avoir aucune
racine commune avec une équation rationnelle, sans la diviser.}

Car le plus grand commun diviseur entre l'équation irréductible
et l'autre équation sera encore rationnel; donc, etc.

\medskip
\noindent\textsc{Lemme II.} --- \emph{Étant donnée une équation quelconque, qui
n'a pas de racines égales, dont les racines sont $a,b,c,\ldots$, on
peut toujours former une fonction $V$ des racines, telle qu'aucune
des valeurs que l'on obtient en permutant dans \GalPrintedError{celle}{W09-PE001} fonction
les racines de toutes manières, ne soit égale à une autre.}

Par exemple, on peut prendre
\[
V=Aa+Bb+Cc+\ldots,
\]
$A,B,C$ étant des nombres entiers convenablement choisis.

\GalSourceErrorMark{W09-SE001}
\medskip
\noindent\textsc{Lemme III.} --- \emph{La fonction $V$ étant choisie comme il est indiqué
dans l'article précédent, elle jouira de cette propriété,
que toutes les racines de l'équation proposée s'exprimeront rationnellement
en fonction de $V$.}

En effet, soit
\[
V=\varphi(a,b,c,d,\ldots),
\]
ou bien
\[
V-\varphi(a,b,c,d,\ldots)=0.
\]
Multiplions entre elles toutes les équations semblables, que l'on
obtient en permutant dans celles-ci toutes les lettres, la première
seulement restant fixe; il viendra une expression suivante:
\[
\begin{split}
&[V-\varphi(a,b,c,d,\ldots)]
[V-\varphi(a,c,b,d,\ldots)]\\
&\qquad [V-\varphi(a,b,d,c,\ldots)]\ldots,
\end{split}
\]
symétrique en $b,c,d,\ldots$, laquelle pourra, par conséquent,
s'écrire en fonction de $a$. Nous aurons donc une équation de la
forme
\[
F(V,a)=0.
\]
